\clearpage
\fancyhead[L]{\small STEINHEIL. PRISM PHOTOMETER.}
\SourceLeaf{654}
\begin{center}\small Göttingische gelehrte Anzeigen. Nos. 34. 35. pp. 329..337. 5 March 1835.\\[3pt]\rule{.72\linewidth}{.3pt}\end{center}

Göttingen. Three entries were received in answer to the principal prize question set for November 1834 by the Mathematical Class of the Royal Society, although the deadline had been extended until the end of December according to No.~149 of these Notices for the preceding year. One was in Latin, with the motto: Opinionum commenta delet dies, naturae iudicia confirmat. The second was in German, with the inscription: Suum cuique. The third was likewise in German, with the words: Only impressions of the same kind are comparable.

Paper No.~2, with the inscription Suum cuique, contains only its author's opinions on the formation and natural constitution of the celestial bodies, and nothing at all bearing on the solution of the problem set by the Society. A separate assessment of those opinions is therefore unnecessary, since they have no connection whatever with the prize question.

The author of paper No.~1, Opinionum commenta, etc., on the other hand, understood the question correctly, proposed an apparatus for comparing the light intensities of two stars and had it made, and also reported some trials of its application to actual measurements of light. The instrument is a telescope fitted with devices that allow both stars to be seen side by side in the field at the same time, one directly and the other by reflection. The reflection is produced by a mirror mounted in front of the objective. By rotating the mirror about an axis that intersects the line of sight at right angles, it can be set at the inclination to the line of sight corresponding to the angular separation of the two stars. The outer edge of the mirror coincides with this axis of rotation, so in every position the mirror shades half the mirror from direct light. There is also a semicircular diaphragm mounted immediately in front of the objective. It leaves only half the objective open and can be turned all the way around. The amount of this rotation is measured on a graduated ring (just as the amount of the mirror's rotation is measured on an arc graduated in degrees). When the ring's index is at zero, no direct light enters the telescope, whereas after half a turn no reflected light enters it. At every intermediate position the reflected and direct light share the open half of the objective in proportion to the deviations from those two positions. It is easy to see, then, that once the objective diaphragm has been turned so that both stars appear equally bright, the ratio of the two stars' light intensities can be calculated up to a still unknown factor that depends on the attenuation of the light by reflection. This unknown factor is found or eliminated by using a second observation in which only the stars are interchanged. For certain cases the author has added a second mirror so that one star is seen by double reflection, which otherwise makes no difference to the method. An equatorial mount makes the instrument much easier to use.

It is regrettable that the author received this instrument from its maker too late to carry out a thorough test by making numerous measurements. He determined the light ratios of seven pairs of stars from a total of only 44 observations, which are reported only in summary. The results, taking the star listed first in each case as unity, are as follows.
\SourceLeaf{655}
\begin{center}
\begin{tabular}{lc}
Stars & Light ratio\\[3pt]
Rigel, Procyon & $0.8501$\\
Rigel, $\beta$ Canis Minor & $0.1258$\\
Sirius. Rigel & $0.2875$\\
Sirius, Procyon & $0.2756$\\
Procyon, Regulus & $0.3781$\\
Procyon, North Star & $0.4369$\\
Regulus, North Star & $0.5720$
\end{tabular}
\end{center}

The altitudes of the stars, or the quantities on which they depend, are missing. According to the author's account, the probable errors of these determinations, insofar as they can be established by comparing the individual observations with one another, would vary between $\frac{1}{16}$ and $\frac{1}{64}$ of the whole. Yet comparing the first, third, and fourth determinations with one another shows that much larger corrections are necessary, and the last three determinations cannot be reconciled at all. The author himself admits that he does not know how to explain this contradiction. Although one must hope that, after much more extensive trials, he will in future succeed in finding the source of such errors, the suitability of the apparatus for measuring the brightness of luminous points remains unestablished for the present.

The author of the third paper, with the motto \emph{Only impressions of the same kind are comparable}, proposed and built two quite different instruments. He calls one the eyepiece apparatus and the other the prism photometer. Although both can be used for the stated purpose, the former is really intended less for comparing the light intensities of luminous points than for comparing the specific brightness of more extended surfaces, such as the sky background. It will therefore suffice to indicate here only the main features of the second apparatus. The basic idea of this instrument is the familiar observation that a star which appears as an indivisible luminous point to the unaided eye, or through a telescope set for clear vision, spreads into a circular image when the eyepiece is given a position other than that required for clear vision. The farther the eyepiece is from its normal position, the larger the image becomes and, for that very reason, the fainter it is in its parts. For stars of unequal brightness, the eyepiece must therefore be moved to unequal distances from its normal position to make the images appear to have the same surface brightness. In this way the light intensities of two stars can already be compared to some extent by observing blurred images of them \emph{in succession}, making their surface brightness equal as far as the remembered impression permits, and measuring the corresponding eyepiece positions. Naturally one expects little accuracy from so crude a procedure, and is therefore surprised that the trials reported by the author show much closer agreement than one might have expected. This already inclines one to favor the apparatus ingeniously arranged by the author, with which such images of two stars can be seen \emph{simultaneously} and brought to the same surface brightness.

The objective is cut into two equal halves. Unlike those of a heliometer, the halves can be moved independently along their common axis instead of side by side. The midpoint of their travel, which is measured precisely by scales on the outside of the telescope tube, corresponds, when the eyepiece tube is fully inserted, approximately to the position relative to the telescope tube required for a sharp image. The two objective halves receive their light through mirrors whose reflecting surfaces are inclined at $45^\circ$ to the axis of the tube. One of the mirrors (the one farther from the objective) can be turned about this axis by a measurable amount. In observing two stars, this axis must therefore always be directed toward one pole of the great circle joining them. The \SourceLeaf{656}mirrors themselves are glass prisms. The light enters and leaves them perpendicular to their surfaces. Between the objective halves and their respective prism mirrors are diaphragms formed by two pairs of slides. Each pair of slides is moved by a single screw with oppositely cut threads, so that, as the diaphragm forms a larger or smaller right triangle, the midpoint of its hypotenuse remains fixed.

With this arrangement, when the telescope tube and mirrors are suitably positioned and the objective halves depart from their normal position relative to the eyepiece, one sees both stars at once, each as a right-triangular area. Both the apparent size of the triangle and its surface brightness depend on this displacement. The former also depends on the diaphragm aperture, and the latter on the brightness of the star itself. One can therefore equalize the surface brightness of the two images by altering one of the displacements and, if desired, also equalize their size by altering one diaphragm aperture. It needs no further explanation how the ratio of the light intensities of two stars can thus be found and any differences between the objective halves and prism mirrors eliminated by reversing the combination.

The author subjected his apparatus to a rigorous test, deliberately using artificially produced star-like luminous points in place of stars. He obtained these artificial stars by the reflection of daylight from two nearly equal, well-polished steel spheres about $\frac{1}{2}$ inch in diameter. The daylight came from the same part of the sky background for both spheres and reached them through circular stops of different apertures. Care was taken to prevent stray light from reaching either the spheres or the observer's eye. Four stops were used in all, the smallest 7 and the largest 20 lines in diameter. The six different combinations could thus yield artificial stars with six different ratios of light intensity. According to the author's own investigations, the greatest disparity, 1 to 8, is close to the mean ratio for two stars two magnitude classes apart. These artificial stars do indeed appear very similar, but without the fluctuations and wavering that often make observations of actual stars so uncertain. They also had the very important advantage that their brightness ratio was known a priori from the openings in the stops. The author reports in full detail the large number of measurements of their light intensity made with the prism photometer, without suppressing those that at first showed some irregularities, whose causes were, however, discovered and removed. The measurements as a whole give the probable error of a single comparison as $\frac{1}{32}$ of the whole brightness, whether it is great or small, and the ratios for the different artificial stars agree entirely satisfactorily with the openings in the stops.

The apparatus has thus been shown to be suitable for precise comparison of the brightness of luminous points. Though one may regret the lack of applications to actual stars, there is good reason to expect satisfactory results for them as well, provided that, as the author rightly requires, observations are restricted to especially favorable atmospheric conditions. Given the ease of handling the instrument, one will accomplish more in a few hours under such conditions than in many days under unfavorable ones. The paper also contains a number of other photometric investigations and views of considerable interest. As they are inessential to the main subject, they may be passed over in silence here. Some applications of the calculus of probabilities in the final section would require correction, but this is quite immaterial to the main subject itself.

Finally, it may be noted that the prism photometer, although based on a wholly different
\SourceLeaf{657}principle from that underlying paper No.~1, at the same time permits the comparison of stars by the other principle, namely by means of \emph{sharp} images appearing simultaneously with a measurably reduced objective aperture. Even in this mode of observation, which the author, incidentally, considers unacceptable in light of his experience, the design of the prism photometer would have advantages over that described in connection with paper No.~1.

As paper No.~3 solved the problem most completely and did so in a way that lays the foundation for a valuable advance in this branch of practical astronomy, the Royal Society awarded it the prize and granted the accessit to paper No.~1, which is also highly meritorious.

According to the note unsealed at the Society's public session on 14 February, the author of the prize-winning paper is

\begin{center}
Dr.\ \textsc{Steinheil} in Munich.
\end{center}

The note accompanying paper No.~2 was burned unopened at the same session.
